\vspace{-1mm}
\section{Experiments and Results}
\label{sec:experiments}
\label{sec:results}

\subsection{Experimental Settings}
\label{sec:experiments:settings}

\subsubsection{Datasets}
\label{sec:experiments:datasets}

We conduct experiments on two synthetic/phantom datasets, \textbf{Liver} and
\textbf{Spine}, and two real-world datasets, \textbf{Arm} and
\textbf{Breast}. Liver and Spine are adopted from
UltraNeRF~\cite{wysockiUltraNeRFNeuralRadiance2024}: Liver contains 6
CT-based simulated sweeps with 200 frames each at $480\times640$, whereas
Spine contains 8 robotic sweeps of a lumbar-spine phantom with 150 frames
each at $168\times1300$. Arm and Breast are derived from
TUS-REC-2024~\cite{liTUSREC2024Challenge2025} and
TDSC-ABUS-2023~\cite{luoTDSCABUSChallenge2025}, respectively. Arm consists of
freehand forearm scans from 85 volunteers at $480\times640$, while Breast
contains automated 3D breast ultrasound volumes with typical sizes of
approximately $850\times600\times300$ voxels. We use all 6 Liver and 8 Spine
sweeps, together with 48 Arm sweeps and 30 publicly available Breast volumes.
All images are adaptively downsampled to approximately $10^5$ pixels per frame
to balance reconstruction efficiency and GPU memory usage while preserving the
acquisition geometry.

\subsubsection{Baseline Methods}
\label{sec:experiments:baselines}

We compare our method with seven representative baseline methods from three
categories: NeRF-based Methods, Ultrasound Gaussian Methods, and Adapted
Gaussian Methods. For NeRF-based Methods, we include
UltraNeRF~\cite{wysockiUltraNeRFNeuralRadiance2024} and
NeRF-US~\cite{dagliNeRFUSRemovingUltrasound2024}, both of which
combine implicit neural fields with ultrasound-specific rendering. For
Ultrasound Gaussian Methods, we compare with
UltraGauss~\cite{eidUltraGaussUltrafastGaussian2025},
UltraGRay~\cite{duelmerUltraGRayPhysicsBasedGaussian2026}, and
UltraGS~\cite{yangUltraGSRealTimePhysicallyDecoupled2025}, which respectively
emphasize plane-based aggregation, physics-based ray casting, and wave-inspired
Gaussian rendering. For Adapted Gaussian Methods, we adapt vanilla
3DGS~\cite{kerbl3DGaussianSplatting2023} and
GaussianPile~\cite{kongGaussianPileUnifiedSparse2026} to ultrasound
reconstruction by modifying their data interfaces and renderers.
For fair comparison, all methods are trained and evaluated using the same
measured probe poses.

\subsubsection{Evaluation Metrics}
\label{sec:experiments:metrics}

We evaluate reconstruction quality using PSNR, L1, SSIM, LPIPS, DISTS, and
GMSD. Higher PSNR and SSIM indicate better performance, whereas lower L1,
LPIPS, DISTS, and GMSD are preferred. We additionally report training time and
rendering frames per second (FPS) for computational efficiency.

\subsubsection{Implementation Details}
\label{sec:experiments:implementation}

We initialize each scene with $10^6$ Gaussians and optimize it for 30k
iterations using Adam. For acquisition-guided pseudo poses, we set $\rho=1$,
use the first retained frame as $i_0$, obtain $d_0$ from acquisition metadata,
and adopt the predefined canonical scan direction. For local rendering, we use
a support multiplier of $2.5$ and set $\gamma=0.25$. For global rendering, we
use a support multiplier of $2.0$, 3-point Gauss--Legendre quadrature with a
2-mm integration half-window, $\lambda_{\mathrm{ray}}=0.55$, and
$T_{\min}=0.35$, with $\beta=0.8$ for local-global modulation. Gaussian
adaptation is performed every 2k iterations during the active adaptation stage.
Residual-guided densification uses $q_{\mathrm d}=0.98$,
$\tau_{\omega}=10^{-6}$, a per-event candidate budget of approximately $5\%$
of the current Gaussian population, and a $0.6$-mm scale threshold.
Path-aware pruning uses $\tau_{\alpha}=10^{-4}$, $q_{\mathrm p}=0.10$, and
$q_{\mathrm A}=0.80$. Unless otherwise specified, the same hyperparameters are
used across all sequences without per-sequence tuning. All experiments are
conducted on a single NVIDIA RTX 4090-class GPU.

\input{tables/main/tab_comparison_four_datasets}

\vspace{-3mm}
\subsection{Comparison with Baseline Methods}
\label{sec:results:comparison}

We compare \methodname{} with seven baseline methods from the NeRF-based,
Ultrasound Gaussian, and Adapted Gaussian categories on the Liver, Spine, Arm,
and Breast datasets, in terms of PSNR, L1, SSIM, LPIPS, DISTS, and GMSD,
together with training time and rendering FPS, as reported in
Table~\ref{tab:comparison:main}, with the corresponding qualitative results
presented in Fig.~\ref{fig:comparison:all_datasets}.
Overall, \methodname{} achieves the strongest reconstruction-oriented
quantitative performance and favorable qualitative fidelity across all four
datasets, demonstrating robustness across anatomical targets and acquisition
settings.

\subsubsection{Comparison with NeRF-based Methods}

Compared with UltraNeRF and NeRF-US, \methodname{} achieves consistently
higher reconstruction accuracy across all four datasets, as shown in
Table~\ref{tab:comparison:main}. In particular, \methodname{} outperforms the
strongest NeRF-based baseline on Liver and Breast, improving PSNR from
$22.07\!\rightarrow\!29.35$ dB and $22.00\!\rightarrow\!28.48$ dB,
respectively. On Spine and Arm, \methodname{} further raises PSNR from
27.46 to 31.02 dB and from 29.43 to 34.25 dB, while improving SSIM from
0.692 to 0.790 and from 0.836 to 0.890. It also reduces L1 by 42.9\%, 26.3\%,
35.3\%, and 48.9\% on Liver, Spine, Arm, and Breast, respectively, with
consistent improvements in GMSD. Although their ultrasound-specific renderers
avoid conventional NeRF ray sampling, \methodname{} achieves substantially
higher fidelity with practical training and rendering efficiency.
Similar advantages are visible in
Fig.~\ref{fig:comparison:all_datasets}: NeRF-based methods tend to smooth
tissue textures and weaken fine responses, particularly around Liver
boundaries and thin Spine structures, whereas \methodname{} better retains
sharp interfaces, fine echoes, and surrounding textures, clearly highlighting the
benefit of explicit Gaussian modeling over implicit neural representations.

\subsubsection{Comparison with Ultrasound Gaussian Methods}

Among ultrasound-specific Gaussian approaches, \methodname{} achieves the most
consistent reconstruction-oriented performance across all four datasets, as
reported in Table~\ref{tab:comparison:main}. On Liver, \methodname{}
outperforms UltraGS by 1.33 dB in PSNR
($28.02\!\rightarrow\!29.35$ dB) and improves SSIM from 0.594 to 0.655,
while on Spine it surpasses UltraGRay with PSNR increasing from 30.63 to
31.02 dB and L1 decreasing by 12.5\%. On real data, \methodname{} exceeds
UltraGRay on Arm by 0.43 dB in PSNR, reduces L1 from 0.013 to 0.011, and
raises SSIM from 0.832 to 0.890; on Breast, it further outperforms UltraGauss
by 0.72 dB in PSNR while improving SSIM from 0.782 to 0.801. 
Although UltraGRay favors perceptual distances on simulated data and
UltraGauss offers the highest rendering speed, \methodname{} consistently
achieves the best PSNR, L1, and SSIM with the best or near-best GMSD, while
its 240--353 FPS rendering remains well beyond real-time requirements, thereby
prioritizing reconstruction fidelity while retaining ample speed for
interactive ultrasound visualization.
\methodname{} also shows clear qualitative advantages in
Fig.~\ref{fig:comparison:all_datasets}, retaining sharper boundary profiles
and richer surrounding texture in Liver, more continuous thin responses in
Spine, and more faithful layered structures in the real Arm and Breast
examples, further supporting the benefit of jointly modeling finite-slice echo
formation and path-dependent transmission.

\subsubsection{Comparison with Adapted Gaussian Methods}

Compared with vanilla 3DGS and GaussianPile, \methodname{} outperforms both
baselines across all four datasets, as shown in
Table~\ref{tab:comparison:main}. On Liver, \methodname{} outperforms 3DGS by
2.62 dB in PSNR ($26.73\!\rightarrow\!29.35$ dB) and reduces L1 by 25.0\%,
while on Spine it surpasses the stronger GaussianPile result by 2.80 dB and
decreases L1 from 0.020 to 0.014. On real data, \methodname{} improves 3DGS
on Arm from 33.96 to 34.25 dB with an 8.3\% reduction in L1, and raises
Breast PSNR by 1.40 dB ($27.08\!\rightarrow\!28.48$ dB), while requiring
less training time than both adapted baselines. Qualitatively,
Fig.~\ref{fig:comparison:all_datasets} further shows that \methodname{} better
suppresses local artifacts in Liver, retains thin responses in Spine, and
preserves more coherent echo structures in the real examples.
These improvements demonstrate the advantage of our ultrasound-specific
renderer, whose explicit modeling of finite-slice echo formation and
path-dependent acoustic transmission better matches B-mode image formation and
thereby enables \methodname{} to consistently outperform generic Gaussian
adaptations.

\begin{figure*}[!t]
\centering
\includegraphics[width=\textwidth]{figures/results/paper_all_datasets_2d_rendering.pdf}
\caption{Qualitative comparisons between \methodname{} and seven baseline
methods on four ultrasound datasets. Yellow boxes highlight representative
regions with magnified views. \methodname{} consistently preserves finer
structures and local echo details than the baselines.}
\label{fig:comparison:all_datasets}
\vspace{-4mm}
\end{figure*}

\vspace{-2mm}
\subsection{Ablation Studies}
\label{sec:results:ablation}

In this section, we conduct comprehensive ablation studies on the Liver and Arm
datasets, covering the Local-Global Gaussian Renderer (LGR),
Contribution-Aware Gaussian Adaptation (CGA), and Acquisition-Guided
Pseudo-Pose Parameterization (PPP). We first evaluate their joint contributions
to the complete framework and then examine the internal design choices of each
component in greater detail.

\input{tables/main/tab_ablation_joint_components}

\begin{figure}[!t]
\centering
\includegraphics[width=\columnwidth]{figures/results/paper_ablation_joint_visual_compact.pdf}
\caption{Qualitative ablation of the proposed components on Liver and Arm
cases. Yellow boxes highlight matched magnified regions.}
\label{fig:ablation:joint_visual}
\vspace{-6mm}
\end{figure}

\subsubsection{Joint Component Analysis}

Table~\ref{tab:ablation:joint} evaluates the three proposed components using
the configuration without LGR, CGA, or PPP as the common reference. Among the
individual components, LGR provides the clearest reconstruction gain,
increasing Arm PSNR from 32.516 to 32.768~dB and reducing LPIPS from 0.164 to
0.141; on Liver, it also lowers LPIPS by 13.3\% despite a slight PSNR decrease.
CGA leaves reconstruction quality nearly unchanged while regulating Gaussian
allocation, whereas PPP raises Arm PSNR by 0.209~dB and SSIM from 0.864 to
0.866. With LGR enabled, adding CGA changes Arm PSNR only from 32.768 to
32.764~dB, while adding PPP increases it to 33.004~dB with 0.868 SSIM; the
complete model remains comparable at 32.998~dB and 0.868, confirming the
complementary roles of the three components. This behavior is also reflected
in Fig.~\ref{fig:ablation:joint_visual}, where LGR yields the most evident
improvements in boundary definition and local echo structures, while CGA and
PPP preserve the resulting structural fidelity.

\input{tables/main/tab_ablation_pseudo_pose}

\subsubsection{Acquisition-Guided Pseudo-Pose Parameterization}

We examine how reconstruction changes as measured pose information is
progressively reduced in Table~\ref{tab:ablation:pseudo_pose}, moving from
measured poses (MP) to step (PP-S), axis (PP-A), and uniform-spacing pseudo
poses (PP-U). Replacing MP with PP-S preserves Liver reconstruction while
increasing Arm PSNR from 32.750 to 32.817~dB with unchanged SSIM of 0.866,
showing that part of the measured motion information can be discarded. PP-A
achieves the highest Liver PSNR of 29.774~dB but degrades Arm PSNR and LPIPS
to 32.565~dB and 0.155, suggesting weaker robustness to real acquisition
variation. PP-U ultimately raises Arm PSNR to 32.998~dB and SSIM to 0.868
while remaining close to MP on Liver
($29.748\!\rightarrow\!29.755$~dB), showing that acquisition regularity is
sufficient to preserve geometric consistency and image fidelity without
measured-pose dependence.

\input{tables/main/tab_ablation_local_global_renderer}

\subsubsection{Local-Global Gaussian Renderer}

Table~\ref{tab:ablation:renderer} traces the renderer from local echo modeling
(L) to the addition of global transmission (G), followed by ray integration
(RI) with increasing quadrature order. The local-only formulation attains the
highest Liver PSNR of 29.854~dB but gives 32.525~dB PSNR and 0.164 LPIPS on
Arm, exposing its limited ability to capture propagation-dependent appearance.
Introducing G reduces Liver LPIPS from 0.268 to 0.232 while improving Arm PSNR
to 32.660~dB and LPIPS to 0.141, confirming the value of propagation-aware
acoustic modulation. Ray integration further increases Arm PSNR to
32.759~dB with one-point quadrature, and $N_q=3$ reaches 32.764~dB while
retaining 0.141 LPIPS, suggesting that low-order quadrature provides an
effective approximation to causal ray interaction and thereby improves
path-dependent B-mode rendering.

\input{tables/main/tab_ablation_gaussian_adaptation}

\subsubsection{Contribution-Aware Gaussian Adaptation}

The adaptation study in Table~\ref{tab:ablation:adaptation} follows the
Gaussian population as gradient densification (GD) is replaced by
contribution-based densification (CD), then complemented by local pruning (LP)
and global-path protection (GP). CD changes reconstruction quality only
slightly but expands the final population from 732k to 1,050k on Liver and
from 819k to 1,020k on Arm, reflecting more active capacity allocation. LP
removes 28.3\% and 18.6\% of these populations, respectively, while keeping
PSNR and LPIPS comparable. GP subsequently preserves selected path-relevant
candidates, increasing the populations only from 753k to 783k and from 830k
to 844k, so the final design balances contribution-driven capacity allocation
with local redundancy removal and global utility preservation while
maintaining reconstruction quality.


\noindent\textit{Parameter sensitivity.}
\Cref{fig:ablation:parameter_sensitivity} shows that the production setting is
not an isolated optimum. Synthetic and real spacing curves favor opposite sides
of $\rho=1$, motivating a shared default rather than dataset-specific tuning.
Renderer parameters have the larger response range because they directly alter
signal accumulation, whereas the adaptation heatmaps form broad plateaus around
the selected quantiles and budget. Together, the joint, internal, and
sensitivity studies show that the final configuration combines a material
renderer gain with stable Gaussian allocation and a tracking-free pose model.


\begin{figure}[!t]
\centering
\includegraphics[width=\columnwidth]{figures/results/paper_ablation_parameter_sensitivity_single_column.pdf}
\caption{Parameter sensitivity for pseudo-pose spacing, the local-global
renderer, and contribution-aware adaptation. Values are PSNR changes relative
to the default configuration.}
\label{fig:ablation:parameter_sensitivity}
\vspace{-2.5mm}
\end{figure}